\Askhsh{22035}{2}
\begin{ylh}{1}
    \item 9.2. Το Πυθαγόρειο θεώρημα
    \item 10.1. Πολυγωνικά χωρία
    \item 10.3. Εμβαδόν βασικών ευθύγραμμων σχημάτων.
\end{ylh}
\begin{wrapfigure}[10]{r}{2.8cm} % The height and width are estimates, adjust as necessary
    \centering
    \begin{tikzpicture}[scale=0.04]
        % Ορισμός των κορυφών Α, Β, Γ
        \tkzDefPoint(0,80){A} % Σημείο Α
        \tkzDefPoint(60,80){B} % Σημείο Β
        \tkzDefPoint(60,0){G}  % Σημείο Γ
        
        % Υπολογισμός του σημείου Δ
        % Η διαγώνιος ΑΓ έχει μήκος sqrt((60-0)^2 + (0-80)^2) = sqrt(3600 + 6400) = sqrt(10000) = 100.
        % Δεδομένου ότι ΑΔ = ΓΔ και η γωνία Δ είναι 60 μοίρες, το τρίγωνο ΑΔΓ είναι ισόπλευρο.
        % Συνεπώς, ΑΔ = ΓΔ = ΑΓ = 100.
        % Βρίσκουμε τις δύο τομές των κύκλων με κέντρα Α και Γ και ακτίνα 100.
        \tkzInterCC(A,100)(G,100) \tkzGetPoints{D_up}{D_down}
        \tkzLetPoint{D}{D_down} % Επιλέγουμε το κάτω σημείο για το Δ, όπως στο σχήμα
        
        % Σχεδίαση του πολυγώνου ΑΒΓΔ
        \tkzDrawPolygon(A,B,G,D)
        
        % Σχεδίαση της διαγωνίου ΑΓ
        \tkzDrawSegment(A,G)
        
        % Ετικέτες κορυφών
        \tkzLabelPoint[above left](A){A}
        \tkzLabelPoint[above right](B){B}
        \tkzLabelPoint[below right](G){$\Gamma$}
        \tkzLabelPoint[below left](D){$\Delta$}
        
        % Ετικέτες μηκών πλευρών
        \tkzLabelSegment[above](A,B){$60 \text{ m}$}
        \tkzLabelSegment[right](B,G){$80 \text{ m}$}
        
        % Σήμανση ορθής γωνίας στο Β
        \tkzMarkRightAngle(A,B,G)
        
        % Σήμανση γωνίας 60 μοιρών στο Δ
        \tkzMarkAngle[arc=ll,size=1.2cm](G,D,A)
        \tkzLabelAngle[pos=1.5](G,D,A){$60^\circ$}
        
        % Σήμανση ίσων πλευρών ΑΔ και ΓΔ
        \tkzMarkSegments[mark=s|,size=4pt](A,D G,D)
        
    \end{tikzpicture}
\end{wrapfigure}
Το τετράπλευρο ΑΒΓΔ του σχήματος παριστάνει την κάτοψη ενός κτήματος με $AB = 60 \text{ m}$, $B\Gamma = 80 \text{ m}$, $\hat{\Delta} = 60^\circ$ και $\hat{B} = 90^\circ$ και $A\Delta = \Gamma\Delta$.
\begin{alist}
    \item Να υπολογίσετε το μήκος της διαγωνίου ΑΓ. \monades{09}
    \item Να αιτιολογήσετε γιατί το τρίγωνο ΑΔΓ είναι ισόπλευρο. \monades{04}
    \item Να υπολογίσετε το εμβαδόν των τριγώνων ΑΒΓ και ΑΔΓ. Πόσο είναι το συνολικό εμβαδόν του κτήματος; \monades{12}
\end{alist}